\clearpage
\begin{proofbox}
\textbf{配方法证明（正文）}\quad 令
\[
d=\bar y-a-b\bar x.
\]
则
\[
y_i-a-bx_i=\bigl[(y_i-\bar y)-b(x_i-\bar x)\bigr]+d.
\]
因为 $\sum(x_i-\bar x)=0$、$\sum(y_i-\bar y)=0$，平方展开后的交叉项相消：
\[
Q(a,b)=\sum\bigl[(y_i-\bar y)-b(x_i-\bar x)\bigr]^2+nd^2.
\]
固定 $b$ 时，第一项不随 $a$ 改变，而 $nd^2\ge0$；故最优时 $d=0$，即 $a=\bar y-b\bar x$。这也解释了回归线必过样本中心：总误差有一笔来自“直线在中心处太高或太低”，先把这笔降到 $0$。

代入后，记 $S_{yy}=\sum(y_i-\bar y)^2$，得到
\[
Q(b)=S_{yy}-2bS_{xy}+b^2S_{xx}
=S_{xx}\left(b-\frac{S_{xy}}{S_{xx}}\right)^2
+S_{yy}-\frac{S_{xy}^2}{S_{xx}}.
\]
当 $S_{xx}>0$ 时，平方项仅在 $b=S_{xy}/S_{xx}$ 时为 $0$，所以唯一取得最小值。故
\[
\hat b=\frac{S_{xy}}{S_{xx}},\qquad \hat a=\bar y-\hat b\bar x.
\]

\textbf{拓展阅读：偏导数版本}\quad 下图把同一结论用二元函数的观点表示。右图的横轴只标 $c$，其中 $c=a+\bar x b$；在 $(c,b)$ 坐标中，误差碗的等高线满足
\[
Q-Q_{\min}=n(c-\bar y)^2+S_{xx}(b-\hat b)^2.
\]

% R038：拓展阅读的偏导几何作用；等高线使用 c=a+xbar b 与 b 坐标，因而可轴对齐。
\begin{center}\begin{tikzpicture}[x=1cm,y=1.08cm,>=Stealth,every node/.style={font=\small}]
\path[use as bounding box] (0,0) rectangle (12.4,7.0);
\node[font=\large\bfseries,text=CProof] at (6.2,6.65) {拓展：两个偏导分别调截距与倾斜；误差碗只有一个最低点};
\draw[rounded corners=5pt,fill=purple!3,draw=CProof] (.15,.28) rectangle (6.02,6.12);
\node[font=\bfseries,text=CProof] at (3.08,5.72) {对 $a$：固定 $b$，直线上下平移};
\draw[->,black!65] (.70,1.02)--(5.45,1.02);\draw[->,black!65] (.70,1.02)--(.70,4.70);
\foreach \x/\y in {1.20/1.82,2.00/2.32,2.80/2.08,3.60/3.20,4.50/3.56}{\fill[blue!80!black] (\x,\y) circle (2.5pt);}
\draw[gray!65,thick] (.95,1.58)--(5.10,4.03);\draw[CStatement,very thick] (.95,1.35)--(5.10,3.80);\draw[gray!65,thick] (.95,1.12)--(5.10,3.57);
\fill[orange!90!black] (3.05,2.59) circle (3.4pt);\node[above left,text=orange!85!black] at (3.05,2.59) {中心};
\draw[red!78!black,thick] (2.0,2.32)--(2.0,1.97);\draw[red!78!black,thick] (3.6,3.20)--(3.6,2.92);\node[text=red!78!black,font=\scriptsize,anchor=west] at (3.16,3.38) {残差示意};
\draw[rounded corners=5pt,fill=purple!3,draw=CProof] (6.40,.28) rectangle (12.25,6.12);
\node[font=\bfseries,text=CProof] at (9.32,5.72) {$c=a+\bar x b$ 与 $b$ 的误差碗};
\draw[->,black!65] (6.98,1.05)--(11.68,1.05) node[right] {$c$};\draw[->,black!65] (6.98,1.05)--(6.98,4.82) node[above] {$b$};
\draw[CProof] (9.35,2.92) ellipse [x radius=.78,y radius=.42];\draw[CProof] (9.35,2.92) ellipse [x radius=1.42,y radius=.80];\draw[CProof] (9.35,2.92) ellipse [x radius=2.02,y radius=1.18];
\fill[orange!90!black] (9.35,2.92) circle (3.7pt);\node[above right,text=orange!85!black] at (9.35,2.92) {$(\bar y,\hat b)$};
\draw[->,red!75!black,very thick] (11.16,4.18)--(10.02,3.16) node[midway,above right,font=\scriptsize] {误差下降方向};\draw[->,red!75!black,very thick] (7.55,1.55)--(8.70,2.58);
\end{tikzpicture}\end{center}


令 $e_i=y_i-a-bx_i$。由 $\partial Q/\partial a=0$ 与 $\partial Q/\partial b=0$，分别得
\[
 \sum_{i=1}^n e_i=0,\qquad \sum_{i=1}^n x_i e_i=0.
\]
第一式化为 $n\bar y-na-bn\bar x=0$，即 $a=\bar y-b\bar x$。把它代入第二式，或先将第二式减去 $\bar x\sum e_i=0$，得到
\[
 \sum_{i=1}^n(x_i-\bar x)\bigl[(y_i-\bar y)-b(x_i-\bar x)\bigr]=0.
\]
因而
\[
 b\sum_{i=1}^n(x_i-\bar x)^2=\sum_{i=1}^n(x_i-\bar x)(y_i-\bar y).
\]
由 $S_{xx}>0$ 可除得 $b=\hat b$，再由 $a=\bar y-b\bar x$ 得 $a=\hat a$。

\clearpage
% R038：复用例1数据 (1,2),(2,3),(3,5),(4,6)，中心为 (2.5,4)。
\begin{center}
\begin{tikzpicture}[x=1cm,y=1cm,>=Stealth,every node/.style={font=\small}]
\path[use as bounding box] (0,0) rectangle (12.4,6.35);
\node[font=\large\bfseries,text=CProof] at (6.2,6.02) {中心化：样本中心移到新原点，样本之间的相对位置不变};

\draw[rounded corners=5pt,fill=purple!2,draw=CProof] (.18,1.25) rectangle (5.82,5.42);
\node[font=\bfseries,text=CProof] at (3.00,5.04) {原坐标：例1的四个样本点};
\draw[->,black!65] (.82,1.70)--(5.32,1.70) node[right] {$x$};
\draw[->,black!65] (.82,1.70)--(.82,4.95) node[above] {$y$};
\foreach \x/\y in {1.40/2.70,2.40/3.20,3.40/4.20,4.40/4.70}{\fill[blue!80!black] (\x,\y) circle (2.8pt);}
\fill[orange!90!black] (2.90,3.70) circle (4.0pt);
\node[anchor=south,text=orange!85!black,fill=purple!2,inner sep=1pt] at (2.90,3.76) {$(\bar x,\bar y)=(2.5,4)$};

\draw[->,CProof,very thick] (6.03,3.63)--(6.43,3.63) node[midway,above=3pt,font=\scriptsize] {平移};
\node[align=center,text=CProof,font=\scriptsize] at (6.23,2.98) {$(x_i,y_i)\mapsto$\\$(x_i-\bar x,y_i-\bar y)$};

\draw[rounded corners=5pt,fill=purple!2,draw=CProof] (6.58,1.25) rectangle (12.22,5.42);
\node[font=\bfseries,text=CProof] at (9.40,5.04) {中心化后：新坐标轴经过 $O$};
\coordinate (O) at (9.40,3.70);
\draw[->,black!65] (7.18,3.70)--(11.85,3.70) node[right] {$x'$};
\draw[->,black!65] (9.40,1.72)--(9.40,4.95) node[above] {$y'$};
\foreach \x/\y in {7.90/2.70,8.90/3.20,9.90/4.20,10.90/4.70}{\fill[blue!80!black] (\x,\y) circle (2.8pt);}
\draw[orange!90!black,line width=1.1pt] (O) circle (4.0pt);
\node[below left,text=orange!85!black] at (O) {$O=(0,0)$};

\node[align=center,text=CProof] at (6.2,.66) {中心化后 $\sum x_i'=0,\ \sum y_i'=0$；橙色的样本中心准确移到 $O$，\textbf{它不是原始样本点。}};
\end{tikzpicture}
\end{center}


中心化还给出常用等价式：
\[
 \sum(x_i-\bar x)(y_i-\bar y)=\sum x_iy_i-n\bar x\bar y,\qquad
 \sum(x_i-\bar x)^2=\sum x_i^2-n\bar x^2.
\]
将 $\hat a=\bar y-\hat b\bar x$ 改写为 $\bar y=\hat a+\hat b\bar x$，即回归线经过样本中心；因此在这条约束下改变斜率时，整条候选线才是绕样本中心旋转。
\end{proofbox}
